
%% ================================================================
%%  7.  DISCUSSION
%% ================================================================
\section{Discussion}\label{sec:discussion}

Our benchmark reveals a nuanced picture: no single paradigm dominates
across all settings. Rather, the optimal choice of solution method
depends critically on the demand regime, network topology, and
operational constraints. We synthesize the key findings below.

\subsection{Where RL/DL Methods Excel}

Our results identify three conditions under which learning-based
methods provide significant advantages:

\begin{enumerate}
  \item \textbf{Endogenous demand dynamics (goodwill).} When service
        quality feeds back into future demand, RL
        agents---especially the GNN---learn the goodwill dynamics
        implicitly from experience, while LP-based agents (which
        assume exogenous demand) suffer 25--45\% optimality gaps. This
        is arguably the most practically relevant finding, as real-world
        demand often exhibits endogenous features such as customer
        churn, brand loyalty, and word-of-mouth effects.

  \item \textbf{High-volatility non-stationary demand.} Under demand
        shocks, the GNN and Residual~RL agents degrade by 12--15\%,
        whereas MSSP and DLP degrade by 18--30\%. The distributional
        assumptions underlying LP-based methods are violated during
        the shock period, while RL policies trained with domain
        randomization have already encountered similar disruptions.

  \item \textbf{Real-time decision-making.} RL inference is
        100--300$\times$ faster than MSSP solving. For applications
        requiring sub-second decisions (e.g., e-commerce fulfillment,
        automated warehousing), RL is the only practical approach
        among the methods tested.
\end{enumerate}

\subsection{Where Classical OR Methods Remain Superior}

Conversely, classical OR methods dominate in:

\begin{enumerate}
  \item \textbf{Stationary demand with known parameters.} When the
        demand distribution is known and stable, MSSP achieves
        near-optimal performance ($<6\%$ gap). RL agents cannot match
        this without substantially more training data or architectural
        refinement.

  \item \textbf{Simple topologies with matching structural assumptions.}
        On the serial network, the echelon base-stock heuristic---a
        closed-form solution that is theoretically optimal under
        stationarity---often matches or outperforms RL, as the problem
        structure aligns perfectly with the Clark--Scarf
        assumptions.

  \item \textbf{Interpretability and auditability requirements.}
        Classical OR solutions can be explained through their
        mathematical formulations and optimality conditions.
        For regulated industries or high-stakes decisions, the
        interpretability of LP solutions or heuristic rules may be
        preferred over the black-box nature of neural network policies.
\end{enumerate}

\subsection{The Value of Hybrid Approaches}

The Residual~RL agent demonstrates that combining domain knowledge
with learned corrections is a promising middle ground. In
non-goodwill scenarios, Residual~RL closely matches or slightly
improves upon the raw heuristic---the learned residual adds marginal
value where the heuristic is already near-optimal. In goodwill-enabled
scenarios, however, the residual correction provides meaningful
improvement (2--8 percentage points), as the RL component learns to
compensate for the heuristic's inability to model endogenous demand.

This finding suggests a practical recommendation: \emph{start with a
domain-appropriate heuristic and add an RL residual correction,
rather than training an RL agent from scratch}. The hybrid approach
provides a safe default (the heuristic), faster training convergence,
and interpretable corrections.

\subsection{The Information Structure Effect}

The contrast between informed and blind agents quantifies the
practical cost of demand uncertainty. For MSSP, the blind variant
loses 2--5\% profit under stationary demand but 10--15\% under
non-stationary patterns (trend, shock). This suggests that investment
in demand forecasting infrastructure has higher ROI under volatile
demand---a practically relevant insight.

Interestingly, RL agents are inherently ``blind'' (they estimate
demand from observations) yet remain competitive with or superior to
blind OR methods. This is because RL implicitly learns demand
estimation as part of its policy, whereas blind OR methods apply a
generic rolling-average estimator.

\subsection{Limitations}

We acknowledge several limitations of this benchmark:

\begin{enumerate}
  \item \textbf{Fixed cost structure.} All scenarios use the same
        holding costs, backlog penalties, and prices. Results may
        differ under alternative cost ratios (e.g., high holding
        cost relative to backlog penalty).

  \item \textbf{Fixed lead times.} Lead times are deterministic and
        constant. Stochastic or time-varying lead times would
        introduce additional complexity that could alter the relative
        performance of methods.

  \item \textbf{Single-product setting.} The current environment
        handles a single product. Multi-product settings with
        substitution and capacity sharing represent an important
        extension.

  \item \textbf{Demand distribution family.} All demand is Poisson.
        Results may differ under heavy-tailed (e.g., negative binomial)
        or correlated demand.

  \item \textbf{RL training budget.} The RL agents were trained on a
        consumer-grade laptop with modest computational resources. More
        extensive hyperparameter tuning and longer training could
        potentially close the performance gap with OR methods in
        stationary settings.

  \item \textbf{Episode length.} The 30-period horizon is relatively
        short. Longer horizons may differentially affect methods that
        rely on long-term planning versus reactive control.
\end{enumerate}

These limitations are deliberate \emph{controlled variables}: they
define the scope of the current benchmark while leaving clear
directions for extension. Because the environment is fully
configurable, each of these dimensions can be systematically varied
in future studies using the same evaluation protocol and agent
implementations.
